\section{Limitations and Future Work}
\label{sec:limitations}

\textbf{Reactive state resets on topology change.} A network's circuit is rebuilt from
scratch whenever its wiring changes anywhere within it, resetting a capacitor's charge or an
inductor's current to zero at that instant; the memristor is deliberately exempted, since
persistent state is its central pedagogical purpose. Carrying reactive state across a
topology rebuild the same way is planned future work.

\textbf{Discrete presets remain the default, alongside an optional numeric entry.} Every
adjustable quantity is still, by default, one of a small number of discrete presets cycled by
right-click; a shift-right-click additionally opens a sign-style text-entry screen accepting
any value in the same clamped range, though this does not yet extend to purely categorical
choices (diode family, waveform shape) with no single underlying numeric value.

\textbf{AC analysis is small-signal only, with fixed op-amp pole values.} Every source other
than the pinned AC Source is stamped at 0~V regardless of its own redstone state, the
standard small-signal convention but one that does not distinguish ``genuinely off'' from
``treated as an AC ground.'' The two-pole op-amp model's gain and pole frequencies are fixed
constants rather than an editable preset, and the memristor's AC case remains a frozen
resistor with no frequency dependence of its own -- future work building
on~\cite{gater2026smallsignal,gater2025bodeplots}.

\textbf{Solver scalability.} The MNA system for each network is solved by dense Gaussian
elimination every tick, adequate for hand-built circuit sizes but scaling as $O(n^3)$ without
exploiting sparsity; a sparse solver or amortized re-solves would be needed for very large
circuits or a heavily populated multiplayer server.

\textbf{Verification remains informal, and memristor model fidelity is limited to the
original formulation.} The closed-form comparisons of Section~\ref{sec:validation} are not
yet a persisted regression suite, and the tool's pedagogical effectiveness has been assessed
only informally; a pre/post evaluation using an established circuits concept inventory,
comparing a CircuitCraft cohort against a conventional schematic-simulator cohort, is a
natural next step. The mod also implements only the original HP linear-drift memristor model,
known to exhibit boundary pinning under sustained excitation (Section~\ref{sec:experiments}
observes this directly); exposing a windowed variant such as Biolek et al.'s~\cite{biolek2009spice}
as an additional preset is planned.

\textbf{Platform dependency.} The mod is tied to a specific, externally maintained version of
Minecraft and Fabric, both of which continue to evolve, requiring ongoing verification effort
(Section~\ref{sec:validation}) to track future game updates.
